%%%%%%%%%%%%%%%%%%%%%%%%%%%%%%%%%%%%%%%%%%%%%%%%%%%%%%%%%%%%%%
%%%%%%%%%%%%%%%%%% MA-0703 Integración %%%%%%%%%%%%%%%%%%%%%%%
%%%%%%%%%%%%%%%%%%%%%%%%%%%%%%%%%%%%%%%%%%%%%%%%%%%%%%%%%%%%%%
\newpage
\subsection*{MA-0703 Integración}
\addcontentsline{toc}{subsection}{MA-0703 Integración}

\hypertarget{MA0703}{}
\begin{refsection}
\begin{table}[H]
\centering{
\begin{tabular}{|ll|}
\hline
\textbf{Año:} IV  & \textbf{Requisitos:} MA-0725 Análisis Real II \\ 
\textbf{Ciclo:} Optativa  & \textbf{Correquisitos:} Ninguno \\ 
\textbf{Tipo de curso:} Teórico & \textbf{Horas presenciales por semana:} 5 \\ 
\textbf{Créditos:} 5 & \textbf{Horas de trabajo independiente por semana:} 10 \\ \hline
\end{tabular}
}
\end{table}

\subsubsection*{Descripción del curso}

Este curso optativo profundiza en la teoría de la medida y la integración abstractas, generalizando la teoría de Lebesgue a espacios de medida más generales. Constituye un tema fundamental de la formación avanzada de la persona matemática, dada la amplitud de sus aplicaciones en análisis, geometría, probabilidad y áreas afines.

El curso ofrece una introducción rigurosa al marco abstracto de la medida y, posteriormente, se enfoca en ejemplos y aplicaciones donde esta teoría resulta indispensable: desigualdades isoperimétricas, teoremas ergódicos, medidas de Hausdorff, medidas de Haar y transporte óptimo. Se enfatiza la resolución de problemas y la conexión con temas de investigación contemporánea.

Se supone dominio de la medida e integración de Lebesgue en $\R^n$, espacios $L^p$ y resultados básicos de análisis real (como los desarrollados en MA-0705 Análisis Real I y MA-0725 Análisis Real II).

\subsubsection*{Objetivos}

\textbf{General:} Desarrollar competencias en teoría abstracta de la medida e integración y en sus aplicaciones clásicas dentro del análisis matemático.

\textbf{Específicos:} Al finalizar el curso se espera que la persona estudiante sea capaz de:

\begin{enumerate}
\item Demostrar resultados centrales de la teoría de la medida en espacios generales.
\item Aplicar la teoría de la medida en problemas de isoperimetría, ergodicidad, medidas de Hausdorff, medidas de Haar y transporte óptimo.
\item Relacionar las distintas formas en que la teoría de la medida interviene en aplicaciones dentro de la matemática.
\item Formular y desarrollar un proyecto de investigación sobre un tema vinculado con los contenidos del curso.
\item Comunicar argumentos analíticos con rigor, tanto de forma oral como escrita.
\end{enumerate}

\subsubsection*{Contenidos}

\begin{enumerate}

\item \textbf{Repaso de integración de Lebesgue (2 semanas):}
\begin{enumerate}
    \item Medida exterior, conjuntos medibles y medida de Lebesgue.
    \item Integral de Lebesgue, teorema de Fubini y teoremas de diferenciación.
\end{enumerate}

\item \textbf{Curvas rectificables e isoperimetría (2 semanas):}
\begin{enumerate}
    \item Desigualdad de Brunn--Minkowski.
    \item Contenido de Minkowski de una curva.
    \item Desigualdad isoperimétrica.
\end{enumerate}

\item \textbf{Teoría abstracta de la medida e integración (3 semanas):}
\begin{enumerate}
    \item Espacios de medida abstracta y teorema de Carathéodory.
    \item Espacios producto y teorema de Fubini en el contexto abstracto.
    \item Continuidad absoluta de medidas y teorema de Radon--Nikodym.
\end{enumerate}

\item \textbf{Teoremas ergódicos (2 semanas):}
\begin{enumerate}
    \item Teorema ergódico en media.
    \item Teorema ergódico maximal y teorema ergódico puntual.
\end{enumerate}

\item \textbf{Medidas de Hausdorff (2 semanas):}
\begin{enumerate}
    \item Definición y propiedades básicas.
    \item Dimensión de Hausdorff.
    \item Curvas de Peano y ejemplos geométricos.
\end{enumerate}

\item \textbf{Medidas de Haar (2 semanas):}
\begin{enumerate}
    \item Medidas de Radon.
    \item Medidas de Haar: existencia y unicidad.
\end{enumerate}

\item \textbf{Transporte óptimo (2 semanas):}
\begin{enumerate}
    \item Problema de Kantorovich--Monge.
    \item Dualidad de Kantorovich y aplicaciones introductorias.
\end{enumerate}

\end{enumerate}

\subsubsection*{Metodología}

\noindent \textbf{Lineamientos metodológicos:} La persona docente a cargo de este curso debe respetar las siguientes pautas:

\begin{enumerate}
    \item El curso combina \textbf{clases magistrales} con \textbf{sesiones semanales de ejercicios}; en estas últimas la persona estudiante presenta y discute soluciones de tareas asignadas con anticipación.
    \item Se espera participación activa: planteamiento de preguntas, discusión de ideas y comunicación rigurosa de argumentos en pizarra o equivalente.
    \item El trabajo independiente en casa es fundamental para la resolución de problemas y la asimilación de la teoría abstracta.
    \item Se incluye un \textbf{proyecto} desarrollado a lo largo del curso, preferiblemente en grupo, con acompañamiento docente, sobre un problema que utilice conceptos y técnicas del programa.
    \item Pueden asignarse lecturas complementarias (sin evaluación sumativa obligatoria) para contextualizar el desarrollo histórico y las aplicaciones de la teoría de la medida.
    \item Cuando las autoridades sanitarias o institucionales impongan restricciones, las sesiones presenciales pueden trasladarse a modalidad virtual según indique la persona docente.
\end{enumerate}

\textbf{Sugerencias metodológicas:} Adicionalmente, se tienen las siguientes recomendaciones para la persona docente a cargo de este curso:

\begin{enumerate}
    \item Explicitar las conexiones entre los contenidos de MA-0705 y MA-0725 y el paso al marco abstracto.
    \item Alternar demostraciones formales con ejemplos que muestren la relevancia de cada bloque temático.
    \item Utilizar plataformas institucionales de mediación virtual para entrega de tareas y materiales, cuando corresponda.
    \item Promover estrategias de evaluación formativa y sumativa coherentes con el marco pedagógico de la carrera: tareas periódicas, presentaciones orales, exámenes parciales y proyecto final.
    \item Fomentar la consulta de referencias clásicas como \cite{Fol99}, \cite{SS05}, \cite{Lie01} y \cite{Vil03}.
\end{enumerate}

\nocite{Bog20}
\nocite{BBT08}
\nocite{Fol99}
\nocite{Lie01}
\nocite{RF10}
\nocite{SS05}
\nocite{Vil03}

\printbibliography[heading=subbibliography,section=\the\value{refsection}]
\end{refsection}
